%==============================================================================
%  Delegated Payments for Repetitive Purchases
%  Single-column IEEE style.
%
%  Compile:  pdflatex report.tex   (twice, for cross-references)
%
%  Rupee sign: pdflatex has no glyph for U+20B9, so \rs prints "Rs." If you
%  compile with XeLaTeX or LuaLaTeX and a font that has it, redefine \rs as
%  \newcommand{\rs}{\textrupee\,} and load fontspec.
%==============================================================================
\documentclass[journal,onecolumn,11pt]{IEEEtran}

\usepackage[utf8]{inputenc}
\usepackage[T1]{fontenc}
\usepackage{geometry}
\geometry{a4paper,margin=1in}
\usepackage{booktabs}
\usepackage{array}
\usepackage{enumitem}
\usepackage{xcolor}
\usepackage{url}
\usepackage{graphicx}
\usepackage{tikz}
\usetikzlibrary{arrows.meta}
\usepackage[hidelinks]{hyperref}

% palette for the sequence diagram
\definecolor{lnUser}{HTML}{3355C8}
\definecolor{lnAgent}{HTML}{6A3FD4}
\definecolor{lnCard}{HTML}{C8442F}
\definecolor{lnMerch}{HTML}{0A7686}
\definecolor{lnHot}{HTML}{FBEAE6}
\definecolor{lnSweep}{HTML}{237A4C}
\definecolor{lnGrey}{HTML}{7A8087}
\definecolor{lnCtl}{HTML}{0D5C66}
\definecolor{lnNet}{HTML}{8A5A10}

\newcommand{\rs}{Rs.\,}

% Confidence markers, per the project's evidence discipline.
\newcommand{\primary}{{\footnotesize\textsc{\textcolor{teal!70!black}{[primary]}}}}
\newcommand{\secondary}{{\footnotesize\textsc{\textcolor{gray!80!black}{[secondary]}}}}
\newcommand{\unver}{{\footnotesize\textsc{\textcolor{orange!80!black}{[unverified]}}}}
\newcommand{\contested}{{\footnotesize\textsc{\textcolor{red!70!black}{[contested]}}}}

\begin{document}

\title{Delegated Payments for Repetitive Purchases:\\
A Control and Evidence Layer for Machine-Initiated Transactions}

\author{Ashutosh Anand\\
\textit{Indian Institute of Technology Kanpur}\\
Supervisor: Prof.\ Vimal Kumar\\
\texttt{support@adsys.in}}

\markboth{Technical and Business Report}{Anand: Delegated Payments for Repetitive Purchases}
\maketitle

\begin{abstract}
Machines that reorder consumables on their own must pay without a human present.
Every payment rail now shipping for autonomous agents enforces \emph{how much} an
agent may spend; none verifies \emph{what it actually bought}. We show that this
distinction is not academic: under a \rs 2{,}000 mandate, a \rs 600 charge against
a basket the owner approved at \rs 118 passes every control in production today,
violating no rule and mapping to no dispute code. We describe a control and
evidence layer that sits above the rail rather than replacing it, comprising a
deterministic order check executed before settlement and a cross-rail receipt
binding the mandate, the approved order and both payment legs. The layer is
implemented over two rails with materially different properties---UPI ReservePay,
which is pre-funded and carries no agent identity, and card Agentic Tokens, which
are post-paid and network-scoped---and we report four empirical findings from that
implementation, including the result that cryptographic signatures on an order do
not prevent overcharging under any signer arrangement. We conclude with the
commercial position implied by the architecture.
\end{abstract}

\begin{IEEEkeywords}
Agentic payments, delegated authorisation, UPI, Unified Payments Interface,
ReservePay, agentic tokens, payment mandates, dispute evidence, machine commerce.
\end{IEEEkeywords}

\section{Introduction}
\IEEEPARstart{T}{he} rails for agent-initiated payment arrived during 2026. Mastercard
open-sourced Verifiable Intent in March~\cite{mc-vi}, establishing cryptographic proof
that a human authorised an agent. Pine Labs launched P3P in June~\cite{p3p}, enabling
agents to complete UPI payments after a single upfront approval. Google's AP2~\cite{ap2},
Stripe and Tempo's MPP~\cite{mpp} and Coinbase's x402~\cite{x402} all converge on the
same HTTP~402 challenge--response pattern.

Each of these answers the question \emph{may this agent spend?} None answers
\emph{did this agent buy what it was told to buy?}

That gap is the subject of this report. Section~\ref{sec:scope} defines the scope
narrowly---repetitive purchases by machines, beginning at the moment of payment.
Sections~\ref{sec:pain}--\ref{sec:use} establish the problem, the field evidence that
it is real, and where it occurs. Sections~\ref{sec:product}--\ref{sec:sec} describe the
product, its architecture, lifecycle, implementation and security properties.
Section~\ref{sec:findings} reports what building it revealed, and
Section~\ref{sec:related} places it against contemporaneous systems.
Sections~\ref{sec:biz}--\ref{sec:risk} cover the commercial model, the programme
context and open risks.

\section{Scope}\label{sec:scope}
This work concerns \textbf{repetitive purchases}: a machine detects that it needs
something and buys it again. Filters, electricity, reagents, water, consumables. It is
not concerned with discovery, price comparison or supplier selection.

The restriction is what makes the problem tractable. An owner cannot approve each
purchase, but can approve the \emph{pattern} in advance---this item, up to this much,
from these suppliers, this often. That advance approval is the object against which a
payment is subsequently checked.

Accordingly, the work begins when the basket is complete and the agent requests
payment, and ends when the money has moved and the evidence is recorded.

\section{Problem}\label{sec:pain}

\subsection{The device manufacturer}
To let a product transact, a manufacturer requires a payment licence, a mandate
system, spend controls a customer will trust, and a dispute path---in every market it
sells into. Manufacturers will not become payment companies in order to sell water
filters. This is the party that will pay for a solution---a claim supported by
fieldwork rather than assertion, in Section~\ref{sec:field}.

\subsection{The owner}
The owner will not permit a machine to spend without visible limits and a means of
revocation. More seriously, the owner has no recourse when the machine buys the wrong
thing: NPCI's UDIR dispute system carries reason codes for failed, unauthorised and
fraudulent transactions, and \textbf{none for an agent's mistake}~\secondary.

The owner also cannot see across rails. A UPI block is visible in the bank
application; a card token in the issuer's; the purchase history in neither.

\subsection{The rail}
UPI has no field for an agent. A ReservePay block authorises \emph{a merchant to
debit}, not \emph{an agent to spend}, and the rail cannot distinguish an
agent-initiated debit from any other~\primary. Authentication occurs once, at block
creation; every subsequent debit is unauthenticated by the human~\primary. The
framework is set out in NPCI Operating Circular~228~\cite{npci-228}.

\subsection{The merchant}
As of mid-2026 no card network has published a binding chargeback rule for a dispute
initiated by an agent, and the one specification that addresses the question in writing
assigns the loss to the merchant~\secondary.

\section{Field Validation}\label{sec:field}
Section~\ref{sec:pain} argues that the device manufacturer is the party that pays.
That claim is not left as an assertion. During a ten-day industry immersion the
following was obtained directly from operators and manufacturers~\primary.

\subsection{Liability lands on the manufacturer, unprompted}
The founder of \textbf{Innexia} stated, without being asked and in substantially the
terms of this thesis, that when an automated purchase goes wrong \emph{the customer
comes after us, not the payment provider}.

This is the load-bearing commercial claim of the entire architecture, volunteered by a
manufacturer rather than argued by us. It establishes both that the pain is real and
that it is felt by the party with a budget.

\subsection{Integration intent}
\textbf{zunpulse} and \textbf{NeoSapien} each stated an intent to integrate.

\subsection{Willingness is gated on control, not on payments}
Across \textbf{15 operators} surveyed, willingness to allow autonomous purchasing was
gated \emph{entirely} on the presence of control and revocation, not on payment
capability, cost, or rail.

\subsection{What this establishes}
The market's objection is not \emph{can the machine pay}---that problem is solved by
several vendors. The objection is \emph{can I stop it, and who pays when it is wrong}.
That is precisely the layer described in Sections~\ref{sec:arch}--\ref{sec:feat}, and
it is the reason a control layer is a product rather than a feature.

\section{Use Cases}\label{sec:use}

\subsection{Selection criteria}
A deployment is a candidate only if all three conditions hold.

\begin{enumerate}[leftmargin=*]
\item \textbf{The amount varies.} If the machine always pays \rs 499, a spend cap
already suffices. The check earns its place only where price is not known in advance.
\item \textbf{No human is present.} If a person confirms each purchase, the person is
the check.
\item \textbf{Error is expensive.} A \rs 40 mistake is a nuisance; a \rs 4{,}000
mistake is a support ticket, a refund, and a disabled feature.
\end{enumerate}

This criterion excludes fixed-price subscription reorder, where nothing varies and
therefore nothing requires checking.

\subsection{Segments}
Table~\ref{tab:seg} groups candidate deployments by whether the machine already
transacts in India.

\begin{table}[t]
\centering
\caption{Candidate deployments by present-day maturity in India}
\label{tab:seg}
\begin{tabular}{@{}p{0.24\linewidth}p{0.70\linewidth}@{}}
\toprule
\textbf{Segment} & \textbf{Characteristics} \\
\midrule
\multicolumn{2}{@{}l}{\textit{Transacting today}}\\
EV charging & Price varies by station and time; already an eligible ReservePay merchant category~\secondary \\
Battery swapping & Two- and three-wheeler networks; high frequency, small tickets \\
Water tanker delivery & Several suppliers, prices move daily, owner typically absent \\
Diagnostic analysers & Reagents and consumables; high value, a wrong order halts a laboratory \\
\midrule
\multicolumn{2}{@{}l}{\textit{Machine deployed, payment not automated}}\\
Water purifiers & Cartridge replacement by usage rather than calendar \\
Vending, micro-markets & Machine reorders its own stock \\
Fleet fuel and tolls & Vehicles transacting against a fleet mandate \\
Industrial consumables & Tooling and spares ordered from telemetry \\
LPG and generator fuel & Refill booked before exhaustion \\
Agricultural IoT & Inputs ordered within an agronomic window \\
\midrule
\multicolumn{2}{@{}l}{\textit{Installed base not yet present}}\\
Smart refrigerators & Groceries by stock level \\
Voice assistants & Household reorder on request \\
Washing machines & Detergent by cycle count \\
\bottomrule
\end{tabular}
\end{table}

Segment sizing has not been performed and no figures are offered~\unver.

\section{Product Definition}\label{sec:product}
The system is an \textbf{authorisation and evidence layer situated between an agent and
a payment rail}. It verifies every machine-initiated payment against what the owner
approved in advance, and produces the record a subsequent dispute requires.

Stated for a device manufacturer: the product allows a device to purchase autonomously
without the manufacturer becoming a payment company.

\subsection{Functional components}
Six components. They constitute a single product because none is sufficient alone.

\begin{enumerate}[leftmargin=*]
\item \textbf{Agent registration.} A cryptographic identity chained to a verified owner
and to the host device. A device may host several agents; revoking one does not affect
the others.
\item \textbf{Authority creation.} Per-purchase ceiling, period budget, category scope
and expiry, written to a ReservePay block on UPI or an Agentic Token on cards from a
single interface. The owner authenticates once.
\item \textbf{The order check.} Two arithmetic tests executed before settlement, as
specified in Section~\ref{sec:feat}.
\item \textbf{The cross-rail receipt.} One signed record binding the mandate, the
approved order and both payment legs, with the amount required to agree across all of
them.
\item \textbf{Owner control surface.} Purchase history presented beside the order each
purchase was made against, and a revocation that fans out across both rails and any
resting float.
\item \textbf{Reconciliation.} Three outcomes---failed, succeeded, unknown---with
refunds releasing the mandate headroom they consumed.
\end{enumerate}

\subsection{Integration surface}
A manufacturer integrates seven endpoints.

\begin{verbatim}
POST /agents                register an agent on a device
POST /authorities           owner grants limits
POST /orders                agent submits the basket
POST /payments              agent requests payment
GET  /agents/{id}/history   for the owner application
POST /agents/{id}/revoke    three-action revocation
GET  /receipts/{id}         dispute evidence
\end{verbatim}

\subsection{Explicit non-goals}
The system is not a payment rail; it never moves money on its own account. It is not a
wallet; on the card path it holds no funds at any point. It is not a protocol; it
consumes HTTP~402, AP2 mandate structures and Verifiable Intent rather than proposing
alternatives. It is not an agent; it does not determine what to purchase.

\subsection{Boundary}
The system engages when the basket is complete and the agent requests payment, and
disengages when settlement has occurred and evidence is recorded. Discovery, comparison
and supplier selection are upstream concerns belonging to whoever authored the agent.

\subsection{Why the components are inseparable}
The order check is worthless without the authority, since there would be nothing to
check against. The authority is of limited value without the check---which is precisely
the gap in every system presently deployed. Both are diminished without the receipt,
since a correctly enforced payment can still lose a dispute for want of evidence.

This is the argument for the components constituting a purchasable product rather than
three separable features.

\section{Architecture}\label{sec:arch}
Figure~\ref{fig:seq} traces a purchase funded over UPI; Figure~\ref{fig:seqcard}
traces the same purchase funded by card. They are not the same flow on two pipes---the
UPI path is pre-funded and holds money briefly with the agent, while the card path
routes a token and holds nothing.


\begin{figure}[t]
\centering
\resizebox{\linewidth}{!}{%
\begin{tikzpicture}[x=1.12cm,y=1.05cm,font=\small,
  >={Stealth[length=2.2mm]},
  lane/.style={draw,rounded corners=2pt,inner xsep=6pt,inner ysep=4pt,font=\small\bfseries},
  life/.style={densely dashed,lnGrey!60},
  msg/.style={->,thick},
  cap/.style={font=\scriptsize,text=black},
  sub/.style={font=\scriptsize\ttfamily,text=lnGrey}]

% exposure band, drawn first so everything sits on top
\fill[lnHot] (-1.35,-8.05) rectangle (13.5,-5.15);
\node[sub,anchor=east,text=lnCard] at (13.4,-5.45) {money is with the agent};

% lifelines
\node[lane,draw=lnUser,  text=lnUser]  (U) at (0,0)  {User};
\node[lane,draw=lnAgent, text=lnAgent] (A) at (4,0)  {Agent};
\node[lane,draw=lnCard,  text=lnCard]  (C) at (8.4,0){Virtual card};
\node[lane,draw=lnMerch, text=lnMerch] (M) at (12.8,0){Merchant};
\foreach \x in {0,4,8.4,12.8} \draw[life] (\x,-0.42) -- (\x,-11.3);

% 1
\node[sub] at (-1.15,-1.15) {1};
\draw[msg,lnUser] (0,-1.15) -- (4,-1.15);
\node[cap,above]  at (2,-1.1)  {grants mandate};
\node[sub,below]  at (2,-1.22) {UPI PIN $\cdot$ the only human step};

% 2
\node[sub] at (-1.15,-2.35) {2};
\draw[msg,lnAgent] (4,-2.35) -- (8.4,-2.35);
\node[cap,above]   at (6.2,-2.3)  {card issued, bound to mandate};
\node[sub,below]   at (6.2,-2.42) {limits and expiry live in the credential};

% 3 self
\node[sub] at (-1.15,-3.55) {3};
\draw[msg,lnAgent] (4,-3.35) -- (5.0,-3.35) -- (5.0,-3.8) -- (4,-3.8);
\node[cap,anchor=west]  at (5.2,-3.4)  {discovers, compares, builds the cart};
\node[sub,anchor=west]  at (5.2,-3.78) {local model $\cdot$ no money involved yet};

% 4 self
\node[sub] at (-1.15,-4.75) {4};
\draw[msg,lnAgent] (4,-4.55) -- (5.0,-4.55) -- (5.0,-5.0) -- (4,-5.0);
\node[cap,anchor=west]  at (5.2,-4.6)  {signs an intent for the finished cart};
\node[sub,anchor=west]  at (5.2,-4.98) {\rs 1,240 $\cdot$ 9 items $\cdot$ merchant named};

% 5
\node[sub] at (-1.15,-5.95) {5};
\draw[msg,lnAgent] (4,-5.95) -- (0,-5.95);
\node[cap,above] at (2,-5.9) {draws \rs 1,240 against the mandate};

% 6
\node[sub] at (-1.15,-6.85) {6};
\draw[msg,lnCard,very thick] (0,-6.85) -- (8.4,-6.85);
\node[cap,above] at (4.2,-6.8)  {funds load onto the card};
\node[sub,below] at (4.2,-6.92) {UPI rail $\cdot$ mandate consumed here};

% 7
\node[sub] at (-1.15,-7.75) {7};
\draw[msg,lnCard,very thick] (8.4,-7.75) -- (12.8,-7.75);
\node[cap,above] at (10.6,-7.7)  {card pays the merchant};
\node[sub,below] at (10.6,-7.82) {card rail $\cdot$ Agentic Token + Verifiable Intent};

% 8
\node[sub] at (-1.15,-8.75) {8};
\draw[msg,lnMerch] (12.8,-8.75) -- (4,-8.75);
\node[cap,above] at (8.4,-8.7) {order confirmed};

% 9
\node[sub] at (-1.15,-9.75) {9};
\draw[msg,lnSweep,densely dashed] (8.4,-9.75) -- (0,-9.75);
\node[cap,above] at (4.2,-9.7) {anything unspent sweeps back};

\end{tikzpicture}}
\caption{One purchase across both rails. Steps~5--7 are the only interval in which the
funds belong to neither the owner nor the merchant.}
\label{fig:seq}
\end{figure}


\begin{figure}[t]
\centering
\resizebox{\linewidth}{!}{%
\begin{tikzpicture}[x=1.12cm,y=1.05cm,font=\small,
  >={Stealth[length=2.2mm]},
  lane/.style={draw,rounded corners=2pt,inner xsep=5pt,inner ysep=4pt,font=\small\bfseries},
  life/.style={densely dashed,lnGrey!60},
  msg/.style={->,thick},
  cap/.style={font=\scriptsize,text=black},
  sub/.style={font=\scriptsize\ttfamily,text=lnGrey}]

\node[lane,draw=lnUser, text=lnUser] (U) at (0,0)    {User};
\node[lane,draw=lnAgent,text=lnAgent](A) at (3.4,0)  {Agent};
\node[lane,draw=lnCtl,  text=lnCtl]  (K) at (6.9,0)  {Control layer};
\node[lane,draw=lnNet,  text=lnNet]  (N) at (10.4,0) {Issuer + network};
\node[lane,draw=lnMerch,text=lnMerch](M) at (13.8,0) {Merchant};
\foreach \x in {0,3.4,6.9,10.4,13.8} \draw[life] (\x,-0.42) -- (\x,-10.4);

\node[sub] at (-1.15,-1.15) {1};
\draw[msg,lnUser] (0,-1.15) -- (6.8,-1.15);
\node[cap,above] at (3.4,-1.1)  {enrols a card, once};
\node[sub,below] at (3.4,-1.22) {biometric ID\&V $\cdot$ passkey bound};

\node[sub] at (-1.15,-2.2) {2};
\draw[msg,lnCtl] (6.9,-2.2) -- (10.3,-2.2);
\node[cap,above] at (8.6,-2.15)  {requests a token for this agent};

\node[sub] at (-1.15,-3.25) {3};
\draw[msg,lnNet] (10.4,-3.25) -- (3.5,-3.25);
\node[cap,above] at (7,-3.2)  {Agentic Token issued};
\node[sub,below] at (7,-3.32) {MDES $\cdot$ agent id, caps, categories, expiry};

\node[sub] at (-1.15,-4.4) {4};
\draw[msg,lnAgent] (3.4,-4.2) -- (4.3,-4.2) -- (4.3,-4.65) -- (3.4,-4.65);
\node[cap,anchor=west] at (4.5,-4.25)  {builds the cart, signs an intent};
\node[sub,anchor=west] at (4.5,-4.63)  {local model $\cdot$ no money involved yet};

\node[sub] at (-1.15,-5.5) {5};
\draw[msg,lnAgent] (3.4,-5.5) -- (6.8,-5.5);
\node[cap,above] at (5.15,-5.45) {order + payment request};

\node[sub] at (-1.15,-6.6) {6};
\draw[msg,lnCtl,very thick] (6.9,-6.4) -- (7.8,-6.4) -- (7.8,-6.85) -- (6.9,-6.85);
\node[cap,anchor=west] at (8,-6.45) {\textbf{order check} --- deterministic, pre-settlement};
\node[sub,anchor=west] at (8,-6.83) {refuse here and nothing reaches a rail};

\node[sub] at (-1.15,-7.6) {7};
\draw[msg,lnCtl] (6.9,-7.6) -- (13.7,-7.6);
\node[cap,above] at (10.3,-7.55) {token presented at checkout};
\node[sub,below] at (10.3,-7.72) {+ Verifiable Intent credential};

\node[sub] at (-1.15,-8.6) {8};
\draw[msg,lnMerch] (13.8,-8.6) -- (10.5,-8.6);
\node[cap,above] at (12.1,-8.55) {authorisation request};
\node[sub,below] at (12.1,-8.72) {network validates scope};

\node[sub] at (-1.15,-9.6) {9};
\draw[msg,lnNet,very thick] (10.4,-9.6) -- (13.7,-9.6);
\node[cap,above] at (12.05,-9.55) {cardholder pays merchant};
\node[sub,below] at (12.05,-9.72) {funds never rest anywhere};

\end{tikzpicture}}
\caption{The card rail. There is no virtual card and no float: the Agentic Token is a
reference to the owner's card, and money moves once, directly. Compare
Figure~\ref{fig:seq}, where steps~5--7 hold funds with the agent.}
\label{fig:seqcard}
\end{figure}

Two funding rails are offered concurrently; the owner selects one per agent. The
control layer is identical in both cases. Their properties differ substantially
(Table~\ref{tab:rails}).

\begin{table}[t]
\centering
\caption{Comparison of the two funding rails}
\label{tab:rails}
\begin{tabular}{@{}lll@{}}
\toprule
\textbf{Property} & \textbf{UPI} & \textbf{Cards} \\
\midrule
Funding model      & Pre-funded              & Post-paid, token-routed \\
Primitive          & ReservePay block (SBMD) & Agentic Token via MDES \\
Money at rest      & Briefly, on agent card  & Never \\
Agent identity     & Absent from the rail    & Carried in the token \\
Merchant scope     & Enforced by us          & Enforced by the network \\
Dispute path       & No code for wrong purchase & 120-day chargeback \\
Transaction cost   & Zero MDR                & Approximately 2\% \\
Licence implication & PPI, or PA at \rs 15\,cr & None; no funds held \\
\bottomrule
\end{tabular}
\end{table}

\subsection{The defining asymmetry}
On UPI, three components must be supplied: agent identity, scope enforcement, and the
order check. On cards, the network already supplies the first two.

\textbf{The order check is the invariant, and it is the only component absent from
both rails.} This observation determines the product boundary: everything else is
rail-specific scaffolding.

\subsection{Where each protocol operates}\label{sec:protomap}
Nothing in Table~\ref{tab:protomap} is proposed by this work. The contribution is the
row marked \emph{ours}, and the fact that the rest is consumed rather than reinvented.

\begin{table}[t]
\centering
\caption{Protocol placement across the two rails}
\label{tab:protomap}
\begin{tabular}{@{}p{0.20\linewidth}p{0.26\linewidth}p{0.46\linewidth}@{}}
\toprule
\textbf{Step} & \textbf{Standard} & \textbf{What it establishes} \\
\midrule
Owner enrolment & Biometric ID\&V, passkeys & The human behind the agent is verified \\
Authority, UPI & ReservePay / SBMD~\cite{npci-228} & Funds blocked, ceiling and expiry set \\
Authority, card & Agentic Token via MDES & Agent identity, caps and categories, network-enforced \\
Cart record & AP2 mandate shape~\cite{ap2} & Signed, non-repudiable record of the basket \\
Transport & HTTP 402~\cite{p3p,mpp,x402} & Machine-readable payment challenge \\
\textbf{Order check} & \textbf{ours} & \textbf{The payment matches the approved order} \\
Authorisation & Verifiable Intent~\cite{mc-vi} & Transaction linked to the human mandate \\
Settlement & UPI switch / card network & Money moves \\
\textbf{Evidence} & \textbf{ours} & \textbf{Both legs bound to one purchase} \\
\bottomrule
\end{tabular}
\end{table}

\subsection{The sequence, end to end}\label{sec:seq}
Stated once, in order, with the rail-dependent steps marked.

\begin{enumerate}[leftmargin=*,itemsep=2pt]
\item \textbf{Enrolment} (once). The owner verifies themselves and binds an
authenticator. \emph{Both rails.}
\item \textbf{Authority.} On UPI, a ReservePay block with the control layer as merchant.
On cards, an Agentic Token minted by the issuer carrying the agent identity and scope.
\emph{Rail-dependent.}
\item \textbf{Detection.} The agent observes that a consumable is depleted.
\emph{Both rails, no money involved.}
\item \textbf{Cart.} The agent selects a supplier and builds a basket. \emph{Both rails.}
\item \textbf{Intent.} The agent signs the finished cart. \emph{Both rails.}
\item \textbf{Order check.} Deterministic evaluation against the signed authority. A
refusal here reaches no rail and leaves nothing to reverse. \emph{Both rails, and this
is the invariant step.}
\item \textbf{Funding.} On UPI, the block is debited onto the agent's virtual card,
consuming mandate headroom. On cards, nothing happens---there is no funding leg.
\emph{UPI only.}
\item \textbf{Payment.} On UPI, the card pays the merchant. On cards, the token is
presented and the network validates scope before the issuer authorises. \emph{Both
rails, different mechanism.}
\item \textbf{Receipt.} A signed record binds the authority, the order and every leg.
\emph{Both rails.}
\item \textbf{Close.} On UPI, residual float sweeps back and headroom is released if the
purchase failed. On cards, nothing to sweep. \emph{UPI only.}
\end{enumerate}

Steps 1, 3, 4, 5, 6, 8 and 9 are identical across rails. Steps 2, 7 and 10 differ, and
steps 7 and 10 exist only on UPI. \textbf{That is the precise sense in which the control
layer is rail-independent and the plumbing is not.}

\section{Payment Lifecycle}\label{sec:life}
A purchase passes through eleven states. Table~\ref{tab:life} gives each with the party
holding the money, which is the property that matters for risk.

\begin{table}[t]
\centering
\caption{Purchase lifecycle states}
\label{tab:life}
\begin{tabular}{@{}lll@{}}
\toprule
\textbf{State} & \textbf{Money held by} & \textbf{Successors} \\
\midrule
\texttt{ORDER SIGNED} & Owner    & \texttt{CHECKING} \\
\texttt{CHECKING}     & Owner    & \texttt{FUNDING}, \texttt{REFUSED} \\
\texttt{REFUSED}      & Owner    & terminal \\
\texttt{FUNDING}      & Owner    & \texttt{FUNDED}, \texttt{REFUSED}, \texttt{RECONCILING} \\
\texttt{FUNDED}       & \textbf{Operator} & \texttt{PAYING}, \texttt{REVOKED} \\
\texttt{PAYING}       & \textbf{Operator} & \texttt{PAID}, \texttt{REFUNDING}, \texttt{RECONCILING} \\
\texttt{REFUNDING}    & \textbf{Operator} & \texttt{REFUNDED} \\
\texttt{PAID}         & Merchant & \texttt{SETTLED} \\
\texttt{RECONCILING}  & Unknown  & \texttt{SETTLED}, \texttt{REFUNDING}, \texttt{REFUSED} \\
\texttt{REFUNDED}     & Owner    & terminal \\
\texttt{SETTLED}      & Merchant & terminal \\
\bottomrule
\end{tabular}
\end{table}

\subsection{The exposure window}
\texttt{FUNDED} through \texttt{PAID} is the sole interval in which the funds belong to
neither the owner nor the merchant. Every risk introduced by the two-leg architecture
resides within it, and every mitigation reduces to shortening it: load the exact basket
total, settle immediately, sweep on failure.

\subsection{Invariants}
\begin{enumerate}[leftmargin=*]
\item \textbf{The check precedes any movement of money.} \texttt{CHECKING} lies between
\texttt{ORDER SIGNED} and \texttt{FUNDING}. A payment failing the order check never
reaches a rail, so refusal costs nothing and leaves nothing to reverse.
\item \textbf{A timeout never triggers a refund.} \texttt{RECONCILING} exists as a
distinct state so that \emph{unknown} cannot be silently collapsed into \emph{failed};
refunding an unknown outcome pays twice.
\item \textbf{Refunds release headroom.} \texttt{REFUNDING} returns both the funds and
the mandate headroom consumed by the funding leg.
\end{enumerate}

\subsection{Revocation}
\texttt{REVOKED} is reachable from any live state and comprises three operations:
terminate the card token, revoke the funding mandate, and sweep any resting float. The
difficult case is revocation during \texttt{FUNDED}, where the block is dead but the
funds are already on the card and remain spendable---which is why the sweep belongs to
revocation rather than following it.

\subsection{Layering}
These are \emph{purchase} states. Each rail leg internally executes its own transaction
state machine. One purchase therefore spans two transactions, and the cross-rail receipt
exists because nothing else records that they were the same purchase.

\section{Core Features}\label{sec:feat}

\subsection{Agent identity and enrolment}
Each agent receives a cryptographic identity chained to a verified owner and to its
host device. One device may host several agents; revoking one does not affect others.

\subsection{The order check}
Two deterministic tests execute before settlement.

\begin{enumerate}[leftmargin=*]
\item \textbf{Does the payment still match the order?} Amount, payee, freshness, and
internal consistency between line items and stated total.
\item \textbf{Does the order fall within the authority?} Per-payment ceiling, category,
period budget, execution count, validity window.
\end{enumerate}

Both are arithmetic. \textbf{No statistical model participates in the approval path},
because an issuer cannot decline a payment on a probability; it requires a reason
admissible in a dispute file.

\subsection{The cross-rail receipt}
A single signed record naming the mandate, the order, the UPI leg and the card leg,
with the amount required to agree across all four. Neither rail retains this
association. It converts a dispute from an argument into a lookup.

\subsection{Reconciliation and refund}
Three outcomes are distinguished, not two. Where the merchant leg definitively failed,
refund immediately. Where it definitively succeeded, no action. \textbf{Where the
outcome is unknown, reconcile}---never refund on a timeout, or the payment is made
twice.

Refunds must release the mandate headroom they consumed. Absent this, a sequence of
failures silently exhausts the owner's periodic cap while their cash balance appears
correct.

\section{Implementation}\label{sec:stack}

\subsection{Protocols adopted}
No new payment protocol is proposed. The following are consumed as given.

\begin{itemize}[leftmargin=*]
\item \textbf{HTTP 402 challenge--response}, common to P3P, MPP and x402~\cite{p3p,mpp,x402}.
\item \textbf{AP2 mandate structure}~\cite{ap2}: signed Intent and Cart records. The
shape is followed; the encoding remains plain signed JSON until interoperability
requires otherwise.
\item \textbf{Mastercard Verifiable Intent}~\cite{mc-vi} on the card leg, binding an
authorisation to the human mandate behind it.
\item \textbf{Model Context Protocol} on the agent side for tool access. It is not a
payment protocol and is not treated as one.
\end{itemize}

\subsection{Stack}
The control layer is implemented in Python with a double-entry ledger in which all
postings must sum to zero, an append-only event log, Ed25519 signatures over
canonically serialised records, and monetary values represented exclusively as integer
paise. Owner-facing control is a web application authenticated by passkey.

\section{Simulation with a Local Model}\label{sec:sim}
The architecture is exercised end to end in a simulator. This section states what is
genuine and what is stubbed, because the distinction determines what the results are
worth.

\subsection{What is real}
\begin{itemize}[leftmargin=*]
\item \textbf{The ledger.} Double-entry, with every posting required to sum to zero.
Balances are re-derived from entries during audit rather than read from a stored field,
so the audit is a check and not a tautology.
\item \textbf{The order check.} The production evaluator, unmodified.
\item \textbf{Signatures.} Ed25519 over canonically serialised records. Tampering is
detected because the signature genuinely fails, not because a flag is set.
\item \textbf{The state machine.} Illegal transitions raise rather than write an
impossible row.
\item \textbf{Reconciliation.} The three-outcome sweep, including the case where the
outcome is genuinely unknown.
\end{itemize}

\subsection{What is stubbed}
\begin{itemize}[leftmargin=*]
\item \textbf{The rails.} There is no NPCI connection and no MDES provisioning. Response
codes resemble the real sets closely enough to be recognisable; simulator-specific codes
carry a distinguishing prefix so nobody mistakes them for switch responses.
\item \textbf{The merchants.} Local services with adversarial behaviours that can be
switched on individually.
\end{itemize}

\subsection{Why the agent runs on a local model}
Three reasons, in order of importance.

First, \textbf{the agent must be permitted to be wrong}. The claim under test is that
safety derives from the mandate chain rather than from the agent's judgement, so an
agent that defends itself proves nothing---a compromised agent will not defend itself.
Running a local model lets the agent be manipulated on purpose and observed.

Second, \textbf{reproducibility}. A local model removes a network dependency and an
external service whose behaviour changes without notice. A deterministic fallback
planner produces a fixed intent, which is what makes the rest of the pipeline testable
offline; it is not a toy stand-in but the planner one wants in continuous integration.

Third, \textbf{the injection surface is the point}. Merchant-supplied product text
enters the agent's context directly. That is the attack, and it cannot be studied
without a model that actually reads the text.

\subsection{The planner interface}
The agent is reached through a single interface: it receives a standing instruction and
an event, and returns a proposal. It imports no rails and holds no write path to the
store; a test asserts this so that it remains true under later modification.

Any planner satisfying that interface can be substituted---a local model, a hosted
model, or the deterministic fallback---without altering the gates, the ledger, or the
reconciliation logic. \textbf{The agent is the most replaceable component in the system,
by design.} That is what it means for enforcement to sit outside the agent's trust
boundary.

\subsection{What the simulation cannot establish}
It cannot establish latency behaviour under real rail conditions, the behaviour of an
actual issuer's authorisation policy, or whether a bank would accept the evidence bundle
in a dispute. Those require counterparties, not code.

\section{Payment Security}\label{sec:sec}

\subsection{Principle}
Enforcement resides in the credential, not inside the agent. A prompt-injected agent
cannot move its own limit, because the limit was never in its custody.

\subsection{Controls}
\begin{itemize}[leftmargin=*]
\item \textbf{Fail-closed evaluation.} An unverifiable signature, unparseable payload
or mismatched reference all deny. An evaluator that approves when confused converts a
failure into an approval.
\item \textbf{Ed25519 over canonical serialisation.} Deliberately asymmetric: a shared
secret would allow the agent to forge the owner's consent.
\item \textbf{Integer arithmetic throughout.} A payment system that rounds can be
gamed by rounding.
\item \textbf{Replay protection} on payment references, so a retry is not a second debit.
\item \textbf{Float minimisation.} Load the exact basket total, settle immediately,
sweep on failure. The card path holds no float at all.
\item \textbf{Revocation fan-out} across both rails and any resting float, with each
step reported independently.
\end{itemize}

\subsection{Limitation}
\textbf{Prompt injection is not prevented.} It remains an open research problem, and
any claim of prevention is overstated. What the architecture achieves is to make
injection unprofitable: an injected agent still cannot construct a payment that
matches an order the owner approved. Injection handling is described as risk scoring,
never as prevention.

\section{Findings}\label{sec:findings}
Four results emerged from implementation and are asserted by tests rather than
argument. The implementation comprises 69 tests across rails, ledger, reconciliation
and the order check.

\begin{enumerate}[leftmargin=*]
\item \textbf{Signatures do not prevent overcharging.} A self-consistent inflated
order---one whose line items genuinely sum to the inflated total---is approved whether
signed by the agent, by the merchant, or by both. A signature establishes authorship,
not truth. What refuses such an order is a per-payment ceiling close to the expected
basket size.
\item \textbf{Routing through an agent-held card destroys merchant scoping.} On the
funding leg the payee is the card, so the mandate can only name the card. A mandate
restricted to one merchant silently becomes unrestricted. A payment reaches a merchant
that the same mandate refuses when paid directly.
\item \textbf{A failed and refunded purchase consumes the periodic cap permanently}
unless headroom is explicitly released on reversal. The owner is made whole in cash
and left short in authority.
\item \textbf{Revocation is three operations, not one:} terminate the card token,
revoke the funding mandate, and return any resting float. Implementations typically
perform the second only.
\end{enumerate}

\section{Related Work and Positioning}\label{sec:related}
Most contemporaneous systems are complements rather than competitors.

P3P~\cite{p3p} provides agentic payment over UPI ReservePay, with the Grantex layer
supplying agent identity, spending limits and audit trail. Mastercard Agent
Pay~\cite{mc-ap} issues Agentic Tokens carrying agent identity and network-enforced
merchant scope. MPP~\cite{mpp} and x402~\cite{x402} address machine-to-machine payment
for API resources---a related but distinct problem, in which the item purchased is the
HTTP resource itself and therefore requires no separate order record.

The observation that instruction origin goes unverified after mandate creation is not
novel and is named as open work in the literature~\cite{drift}. Describing the gap is
therefore not a contribution; only a mechanism is.

\subsection{The structural argument}
A rail that adjudicates disputes concerning its own payments is a party to those
disputes. If the entity moving the money also determines whether the basket was
correct, it judges itself. Card networks are not merchants for the same reason.
Neutrality here is not a marketing position but the precondition for the evidence
having value.

\section{Commercial Model}\label{sec:biz}

\subsection{Who pays}
The buyer and the user are different parties. The owner experiences the problem; the
device manufacturer absorbs its cost in support load, refunds and disabled features.
The manufacturer is therefore the customer. Agent platforms are a later scale channel;
payment service providers are a strategic white-label channel. Consumers do not pay
for safety mechanisms---they decline products that lack them.

\subsection{Go to market}
The sequence is: one pilot in a segment that transacts today; written confirmation of
the no-licence path from a PSP or sponsor bank; publication of the threat model and
findings while the standards conversation remains open; and only then the platform
channel, which must be earned with direct customers because a platform buys a threat
model and incident data rather than a demonstration.

Entry is through the engineer who deployed the agents and feels the risk first;
signature comes from the party carrying the loss; the security function is the gate
and requires a threat model rather than a demonstration. The category has no budget
line, so an existing one is borrowed---fraud and risk, audit and compliance, or
product warranty.

\subsection{Revenue}
UPI is zero-MDR, so no interchange exists to share on the UPI leg; the business is
priced as software. Structures are given in Table~\ref{tab:rev}. \textbf{The figures
are derived rather than validated}~\unver; testing them is a ninety-day objective.

\begin{table}[t]
\centering
\caption{Revenue structures}
\label{tab:rev}
\begin{tabular}{@{}p{0.26\linewidth}p{0.30\linewidth}p{0.36\linewidth}@{}}
\toprule
\textbf{Line} & \textbf{Basis} & \textbf{Rationale} \\
\midrule
Platform fee & Per active agent per month & Grows as the manufacturer ships devices \\
Decision fee & Per authorisation evaluated & Rail-independent \\
Integration & One-time per manufacturer & Covers enablement, qualifies buyers \\
Enterprise licence & Annual, to PSPs and banks & White-label within their product \\
Interchange share & Card legs only & Upside, never foundation \\
\bottomrule
\end{tabular}
\end{table}

Price is defended against the loss rather than the software, benchmarked as basis
points of controlled spend beside card interchange at 150--200 basis points.

\section{Programme Context}\label{sec:prog}
This work is presented within Fintech Spark, administered through the Startup
Incubation and Innovation Centre at IIT Kanpur and funded by Pine Labs. The programme
included a ten-day industry immersion in Pune, hosted by pi-labs.ai, during which the
fieldwork in Section~\ref{sec:field} was gathered.

\emph{Note on names.} \textbf{pi-labs.ai} (the immersion host) and \textbf{Pine
Labs} (the programme funder, and the vendor of P3P) are separate organisations whose
names collide phonetically. They are kept distinct throughout this document, and should
be kept distinct when it is presented.

Pine Labs is simultaneously the programme's funder and the vendor of P3P, the system
most closely adjacent to this work. Both facts bear on positioning.

The relevant public fact is that Pine Labs has stated it is in discussion with card
networks to extend P3P beyond UPI~\secondary. That is precisely the two-rail problem
addressed here. The appropriate posture is therefore complementary rather than
competitive: P3P moves the money and Grantex verifies the agent, while the order check
and the cross-rail receipt address what neither presently covers.

The question that determines the scale of the contribution is whether P3P already
binds a payment to an approved order. Its public documentation terminates at a
quickstart and does not describe such a binding, but absence from public documentation
is not absence from the protocol~\unver. Resolving this is the highest-priority
verification item in the project.

\section{Risks}\label{sec:risk}
Six risks are material. Two are structural and will not resolve with effort; the
remaining four are open questions with a defined way to close them.

The single most consequential is whether P3P already binds a payment to an approved
order. If it does, the contribution narrows to the cross-rail receipt and the control
surface. If it does not, the order check stands as described. The question is
answerable in one conversation and has not yet been asked.

Table~\ref{tab:risk} states each risk with its evidential status and the action that
would settle it.

\begin{table}[t]
\centering
\caption{Principal risks}
\label{tab:risk}
\begin{tabular}{@{}p{0.30\linewidth}p{0.13\linewidth}p{0.49\linewidth}@{}}
\toprule
\textbf{Risk} & \textbf{Status} & \textbf{Resolution path} \\
\midrule
P3P already performs the order check & \unver & Read the published SDK, or ask directly. Determines the width of the contribution. \\
Agent debits may be non-compliant & \contested & Whether ReservePay agent debits fall inside RBI's e-mandate exemption is unresolved~\cite{rbi-auth}. If not, issuer liability attaches in full. \\
Agent dispute liability unassigned & \secondary & No network has published a binding rule. The card path's recourse is real but untested. \\
Card path requires enablement & \secondary & Issuer, acquirer and network provisioning; no sandbox exists. A partnership constraint, not a capital one. \\
UPI path requires capital & --- & Float requires a PPI licence~\cite{rbi-ppi}; direct redistribution requires a PA licence at \rs 15\,crore net worth. \\
Appliance installed base absent & --- & Mitigated by leading with segments that already transact. \\
\bottomrule
\end{tabular}
\end{table}

\section{Conclusion}
The rails for agent payment have arrived; the safeguards have not. Every deployed
system enforces a spending limit, and a spending limit cannot distinguish a correct
purchase from an incorrect one of the same value. We have described a deterministic
order check and a cross-rail evidence record that close that gap without introducing a
new protocol, a new rail, or a new custodian of funds, and we have reported four
implementation findings that constrain how such a system must be built---most
significantly that cryptographic attestation of an order does not, by itself, prevent
an owner from being overcharged.

\section*{Acknowledgment}
The author thanks Prof.\ Vimal Kumar (IIT Kanpur) for supervision, and SIIC IIT Kanpur
and Pine Labs for the Fintech Spark programme.

\begin{thebibliography}{99}
\bibitem{p3p} Pine Labs, ``P3P --- Pine Labs Payments Protocol for AI Agents,''
2026. [Online]. Available: \url{https://www.pinelabs.com/docs/online-payments/ai/p3p}
\bibitem{mc-vi} Mastercard, ``Verifiable Intent,'' 2026. [Online]. Available:
\url{https://www.mastercard.com/us/en/business/artificial-intelligence/mastercard-agent-pay.html}
\bibitem{mc-ap} Mastercard, ``Agentic token framework: driving trusted AI
transactions,'' 2025. [Online]. Available:
\url{https://www.mastercard.com/global/en/news-and-trends/stories/2025/agentic-commerce-framework.html}
\bibitem{ap2} Google Cloud, ``Announcing the Agent Payments Protocol (AP2),'' Sep.
2025. [Online]. Available:
\url{https://cloud.google.com/blog/products/ai-machine-learning/announcing-agents-to-payments-ap2-protocol}
\bibitem{mpp} Stripe and Tempo, ``Introducing the Machine Payments Protocol,'' Mar.
2026. [Online]. Available: \url{https://stripe.com/blog/machine-payments-protocol}
\bibitem{x402} Coinbase, ``x402: a payments protocol for the internet,'' 2025.
[Online]. Available: \url{https://github.com/coinbase/x402}
\bibitem{npci-228} National Payments Corporation of India, ``Enhancement in UPI Single
Block Multiple Debits (UPI Reserve Pay),'' Operating Circular OC~228, FY~2025--26.
\bibitem{rbi-auth} Reserve Bank of India, ``Authentication Mechanisms for Digital
Payment Transactions Directions,'' 2025.
\bibitem{rbi-ppi} Reserve Bank of India, ``Draft Prepaid Payment Instruments
Directions,'' Apr. 2026.
\bibitem{drift} ``Instruction origin verification in delegated agent payments,''
arXiv:2604.15367, 2026.
\end{thebibliography}

\end{document}
